% 来源：OEBPS/Text/Section0001_split_000.xhtml#sigil_toc_id_4

\BookFront{推荐序1}

钱每人每天都在用，人人都离不开，人人都在追求，甚至做梦都想要。但是如何赚钱、如何存钱、如何用钱、如何投资，父母不曾教过我们，学校也没有教过我们。我们在学校里学了那么多知识，语文、数学、英语，甚至是文言文、高数、核物理等高难度的内容，但在开始工作之后，因为在日常生活中基本用不到这些东西，我们会逐渐忘记它们。但是我们天天花钱、天天赚钱、天天想钱，却一点儿知识储备也没有，一点儿实务操作技能也没有。

缺什么就要补什么。要快速补充钱的知识，而且是易学易懂易用的真知，有大师智慧做支撑，有真实案例为证，有具体做法、具体产品可操作，让你读过这本书之后，就成为家里投资理财第一人。

我推荐《钱：7步创造终身收入》这本书，只有一个理由，这本书是最好的投资理财类实务教科书。最好的，没有之一。

这本书的核心和脉络，就是简单7步投资法。这本书分为七大部分，对应的是托尼·罗宾斯总结出来的简单7步投资法。这7步法是托尼·罗宾斯采访巴菲特等50位投资理财大师，把他们的投资理财精髓融会贯通，总结提炼出来的。

好多朋友一看这本书这么厚，就被吓跑了。其实这本书的作者是美国励志演讲第一人。首先，他的客户层次很高，他担任私人顾问，指导培训过很多名人，包括美国前总统克林顿、美国网坛巨星小威廉姆斯、电影巨星莱昂纳多·迪卡普里奥、企业家Salesforce.com公司创始人马克·贝尼奥夫等。他还担任著名对冲基金都铎投资公司创始人保罗·都铎·琼斯21年的私人顾问。其次，他组织的研讨会也是训练营，规模大，效果让参与者震惊。一场活动有5000人甚至上万人一起交流，即时沟通，可以让与会者周末两天持续集中注意力长达50个小时。奥普拉来之前说她根本不可能听上2个小时，但是12个小时之后，她站在椅子上，对着镜头大喊：“这是我一生中最棒的体验！”

读这本书，不像是阅读，更像是在听托尼·罗宾斯演讲，像是听他讲了两天，约50个小时的投资理财课，这和你追韩剧、看美剧差不多。书里讲的很多故事都是发生在我们身边的普通人的平凡事，那些人原来也是不管钱的事，但读了罗宾斯的书之后，他们开始重视钱的事，用心打理钱的事，结果他们逐步实现了财务安全，甚至财务自由。财务自由，就是你再也不用只为了多赚钱而拼命了，你的投资理财会给你带来一生稳定持续的现金收入，相当于自己给自己每月稳定地发工资，不工作照样有一份稳定的收入，你可以工作，但只是为了自己高兴而自愿工作。想想看，为了钱工作和为了爱工作，你的心态会有多么大的不同。为钱工作，你是钱的奴隶，为爱工作，钱是你的奴隶，这就是这本书给我上的第一课。积累了足够多的钱，你就能实现财务自由，钱就是自由。钱是身外之物，但自由不是，你做投资理财，目的不是积累更多的钱，而是自由——为了财务自由，为了心灵自由，为了灵魂自由。为了自由，好好打理你的钱吧。为了财务自由，好好读一读这本书吧。

这本书是美国人写的，也是写给美国人的，我们中国人读了会有用吗？

货币有国别之分，理财的基本手段和原则却没有太大的差别。而且，学习这本书的投资理财之道，并结合中国的国情，你会发现实践操作会更简单，会让你更赚钱。比如，这本书出现最多的是401(k)，这是美国工作者普遍会参加的雇员退休养老金投资计划。它其实就是雇员出一部分钱，雇主也出一部分钱来投资基金，国家再让出一部分所得税。中国目前没有这种雇员养老金投资计划，所以你读到这些东西时，可以把它们简化成自己存钱、自己定期投资基金就行了。你只看其投资的基本原则，书里涉及的美国的那些法律规定你可以完全忽略。另外，作者讲了很多美国各州征收的所得税。中国各省市不征收地方税，所以碰到这方面的内容，你快速翻过去，大概了解一下就行了。所以中国人读这本书会更轻松，运用这本书讲的投资理财方法也更容易。下面我结合这本书的核心——简单7步投资法，一步一步地来说。

第一步，多存钱。

这本书第一部分讲第一步，其核心就三个字，多存钱。这一点，中国人比美国人强多了。中国人的储蓄率是世界上最高的。但是遗憾的是，现在的“80后”和“90后”的收入比他们的父母高多了，可花钱也多多了，所以他们存的钱反而较少。优秀的传统要发扬光大。要赚钱，先要有钱；要投资，先要有本钱。很多刚工作的人常常陷入一个误区，我钱少，存不下钱，或者也有一些工作一段时间的人情况相反，我钱多，收入高，不用存钱。

这本书却告诉我们，收入高的人，比如那些体育明星、演艺明星，一年的收入比很多普通人工作一辈子所赚到的钱加到一起还要多，可是这些明星有很多人后来都破产了。相反，有很多低收入的普通人，由于坚持存钱，把收入相当大的一部分存起来做稳健投资，结果他们一辈子生活得很安稳，到了晚年还积累了几千万美元的财富，还大笔捐赠做慈善事业。

可是，把收入拿出来相当大一部分存起来，自然会减少你的消费支出，也会让你减少很多享受，平添不少痛苦。那么，怎么样才能多存钱又少遭受痛苦？这本书会告诉你一个诺贝尔经济学奖获得者想出来的巧妙办法，即承诺明天储蓄更多。

存钱了，有本钱了，怎么投资理财，让财富增值更多？投资理财市场是全球最大的市场，投资理财产品的广告到处都是，我要告诉你的只有两个字：小心。小心，千万要小心。这就是第二步。

第二步，少上当。

很多人意识到了投资理财的重要性，购买了各种投资理财产品，并把钱交给那些专业机构的专家来打理。可是事实上，很多投资理财产品和大多数医药保健品一样，你信的只是广告，不是疗效。你的投资是赚钱了，但是赚的钱大部分都进入银行、证券公司、基金公司的口袋里了。但是如果亏了钱，对不起，风险自担，没有一个机构会帮你分担。我在证券公司工作十年，慢慢地明白了，传说中的“十个人炒股，七赔二平一赚”是真的。其实“二平”的那两个人是赚钱的，但是扣除手续费等各种费用之后，他们能持平就很不错了。我在基金公司工作十多年了，结果发现，基金很赚钱，但基民大多数不赚钱，因为大多数基民都是在牛市高位买入，在熊市低位卖出，在长期赚大钱的基金上波段操作却亏了不少钱。有的是自己操作，也有不少是受了只想推销基金的银行理财经理的忽悠。还有一个买基金的大误区是，追逐上一年的业绩冠军基金，这就和在生活中购物时选择去年销量冠军的商品一样，其结果是你亏大了。但二者也有不同，购买上年的销量冠军商品，一般情况下，第二年即使不是冠军也至少还是前十名，可是上一年的基金业绩冠军，经常在下一年跌成了倒数的业绩冠军。尽管这个风水轮流转的事实圈里人都知道，但是基金公司和银行还是会大做广告，推销那些季度和年度的业绩冠军基金，很多外行一听就上当了。隔行如隔山，要进哪一行，就得懂那一行的规矩，唯其如此，才能少上当。

这本书第二部分告诉你9个最大的投资营销谎言，戳穿这些谎言，可以让你少上当。这就像孩子刚刚参加工作一样，走正道最重要，倘若孩子上当受骗走上了歪门斜道，他这辈子就完了。你要是听信那些投资营销谎言，你这辈子也完了，因为你的钱都进了证券公司、经纪人、银行理财经理的口袋里了，你是用自己的钱给别人攒退休养老金。

在我们看破了那些投资营销谎言后，如何走上投资理财的光明大道？这好似你的孩子大学毕业后，对事业和生活做了合理的配置，在达到一定层次之后，这个配置就需要升级，于是他向大师学习，力争出人头地，而且他到了一定年纪之后就会明白什么是比金钱更加重要的财富。这个财富正是这本书投资理财的第五步。别急，我们一步一步来。

第三步，做规划。

我们出门要看地图，确定出门的大概路线。我们的事业也要先定个规划，确定好目标，选择好路径。做投资理财也一样，首先你要定出你的财务目标，说得直白一点，目标就是你要实现财务自由梦想需要多少钱，以及运用你现在的投资理财手段，积累到实现财务自由梦想的足够数量的金钱，需要多少时间，包括需要存多少年，投资多少年。实现财务自由梦想，和养孩子、种树一样，需要时间。一夜暴富，跟快速养大的小鸡一样不可靠，和速生林木一样不结实。这一章是我特别喜欢的一章。记住，弄清楚就是知道具体数字是多少，弄清楚了才会有力量。举一个生活中的例子，我看见很多年轻人一玩手机就是几个小时，很多老年人一看电视就是几个小时，可是这些人，不管老少，到了十字路口为了赶绿灯都会飞奔，只为了节省那几秒或者几十秒。是这些人非常珍惜时间吗？可是他们平时浪费的时间数百倍于等红灯的时间。唯一的区别是，平时看电视、看电脑、看手机时，没有像红绿信号灯这样清楚、巨大的数字显示器。我们都有自己的梦想，但都只停留在梦想阶段，我们没有仔细计算过，实现这样的梦想需要花多少钱。这本书把财务梦想分成五级：财务安全、财务活力、财务独立、财务自由、绝对财务自由。我们都向往美好的生活，但是很多人只是想生活美好，并没有算过它需要你花多少钱、需要你有多高的收入、需要你的投资理财有多高的年收益率。这本书为你提供了五级财务梦想的计算表格，也提供了具体案例，你可以照着算一算。从此你的梦想就不再只是梦想了，你开始了梦想转化为现实的第一步。作者考虑得非常周到，很多年轻的朋友照着一算发现，要实现自己的梦想，按照自己目前的收入水平、收入储蓄比例、投资理财收益增长率，需要的时间为好几十年，这样的梦想太遥远了，可能让你有些灰心丧气，但别放弃！作者给你支了五招，它们可以大大加快你的财富积累速度，让你提前几年甚至几十年实现你的财务梦想。加速财富积累的这五招是，存钱更多，收入更高，税费更少，赚的更多，一举多得。

第四步，做配置。

资产配置，听起来很专业，其实相当于衣服搭配好会更美，食物搭配好会更有利于健康，不同的基金或者股票债券等按不同的比例搭配在一起，会让你更赚钱，而且风险更低。

我在证券投资行业干了20多年，这个行业只关注两件事：一个是选股，一个是选时。不信你听一听那些专业投资者和业余投资者的问题，这些问题只有两类：一类是买什么，一类是何时买。时间长了，我才慢慢地明白，慢慢地接受这个事实，选股很重要，选时很重要，但是选股和选时加起来，也没有资产配置的1/10重要。我研究巴菲特有十多年，我现在特别喜欢他的那句名言：未来唯一确定的事是不确定。我们都称巴菲特是股神，巴菲特却从来不这样看自己，巴菲特在他致股东的信中一次又一次地说，我的工作是配置。不同的大师，对投资有完全不同的观点，但是所有的大师都认可一点，那就是资产配置是决定投资长期业绩的最重要因素。如果说专业投资者对选股和选时还有一定能力和经验的话，业余投资者在选股和选时上则完全处于劣势。业余投资者要和专业投资者在同一个市场上竞争，只能靠资产配置，所以业余投资者必须更加重视资产配置。

但是资产配置是投资行业最专业的东西，学校的教授讲起它来都会非常复杂，投资教科书上也都是复杂的模型和公式，专业机构里的专业投资者搞的模型和公式比教科书上的更加复杂，业余投资者要搞懂这些资产配置模型，至少要上金融投资的研究生课程。所以给业余投资者讲资产配置，是投资理财类图书最大的挑战。资产配置非常重要，不能不讲，但它非常复杂，很难讲明白。这本书是我看到的讲资产配置最简单、最明白的一本书。中国有个成语是狡兔三窟，托尼·罗宾斯也讲了三大水桶（中国人将水比作财，所以我特意翻译成水桶）。

第一个水桶是安全/安心水桶。这个水桶求稳，只做固定收益投资，本金安全，收益固定，投资的目的不是进攻，而是防守。

第二个水桶是风险/成长水桶。这个水桶求进，可以做股票投资、房地产投资等，投资的目标是追求更高的收益增长。但是收益高，风险也高，所以托尼·罗宾斯特意把风险放在成长前面，称其为风险/成长水桶。

只求稳，太无聊，而且会让你永远没有机会实现梦想。只求进，太激进，可能会让你赔个精光，连基本的生活也维持不了。所以最关键的是要平衡配置，稳中求进，攻守兼备。具体的配置比例因人而异，而且随着你个人的情况变化而发展，因此你也要与时俱进，及时调整，但是最关键的一点是，你一定要把钱分到两个水桶里，既要安全，又要成长。

与其他投资教科书最大的不同是，托尼·罗宾斯还提出，你要配置一小部分钱到你的梦想水桶，要从现在开始，不要等以后钱多了才去追求梦想，拿出自己能够承担的一小部分钱追逐梦想，而且一旦有笔意外之财，你可以把它放到梦想水桶中。对于那些在大牛市舍不得卖出，结果错过了机会的投资者，这是个好主意。

那么你在做好了资产配置，确定了要买入的投资理财产品，比如股票、基金后，究竟什么时候买入最合适？这是个有巨大争议的话题，因为未来永远是不确定的，所以永远没有确定的答案。

什么样的组合长期来看既很稳定又很赚钱？其实答案就在第五部分前两章。我个人觉得这两章应该归入第四部分，就是瑞·达利欧推荐的稳健投资组合，这是托尼·罗宾斯的这本书的第一大卖点。前面讲的投资理财方法，普通投资者都能用。但普通人不知道，那些有钱人，住的和你不一样，吃的、穿的和你不一样，而且有钱人买的投资理财产品和你也不一样。那些高端投资理财产品，普通投资者从来没有见过，甚至从来没有听说过。达利欧做的是私募基金，投资业绩很显著，投资门槛也很高——只有资产达到50亿美元的投资者才能进来，而且一次至少要买1亿美元的基金。在托尼·罗宾斯的一再劝说下，达利欧终于答应为普通投资者推出他的资产配置普及版。这就和一个中医大师公布了他一生钻研出来的秘方一样，此举让无数人大受裨益。

你做了良好的资产配置，你的资产增值了很大一笔，这是不是可以说你很有钱了，你完全实现财务梦想了？错。资产不是钱，不信，你看看你账上的股票，你再过几年卖出后，你未必能获得现在的价值。再看看你的房子，按照市场价格计算出的它们的价值没有意义，你只有卖出去了，拿到了钱，你才能花。所以，那些有钱人最明白，资产不是钱，真正每月稳稳赚到手的收入才是钱，才是你真正可以花的钱。如何把你的资产转化成可以花的稳定收入？请看第五步。

第五步，稳收入。

我非常喜欢本书第28章所讲的故事，仅攀上山峰不是登山英雄，既上的了山又下的来的，才是登山英雄。你光挣的多还不行，你还要把这些资产转化成未来长期稳定的收入才行。最稳健的投资就是保险。

这本书第28章推荐大家购买一款保险产品——年金，你可以每月得到稳定的收入，而且不管你活多久，只要你活着，就保证能够得到一份稳定的收入。

这本书第29章推荐你购买现代新型年金，这些年金的收益和股票指数涨幅挂钩，你的本金绝对安全，指数涨了你能分享一部分甚至全部上行收益，指数跌了你一分钱也不会亏，你只有上行收益却没有下行风险。原来这种固定指数年金只有有钱人才能买，现在经过托尼·罗宾斯的努力，普通投资者也能购买。

走到第五步，我们已经存了钱，看破了投资理财营销广告的谎言，自己制订好了投资规划，算出了实现财务梦想需要多少钱，知道了用什么方法存钱和投资，明白了需要多少年才能实现梦想，做了合理的资产配置，并且购买了年金保险产品，把资产转化为未来稳定的收入。有了这样高的退休收入，就足够支撑我们梦想的生活方式，从此我们再也不用为钱工作了，我们财务自由了，我们人生自由了，我们心灵自由了。如果你还不满足，还想在投资理财上更上一层楼，在投资游戏上达到更高的水平，你还需要走第六步。

第六步，学大师。

这本书的作者在第1章中提到，为了写作这本书，他采访了50位顶尖投资理财大师，我起初不相信，他是搞励志演讲培训的，是一个投资圈外人，即使是投资圈里人，要见到这么多大师，可能一辈子都无法实现。我真没有想到，托尼·罗宾斯做到了，而且还和这些大师谈得非常深入。比如他采访的约翰·博格，巴菲特称他为自己的英雄偶像，是为美国投资者贡献最大的人。博格原来以为罗宾斯来采访他，不过是和一个外行随便聊聊，结果他们竟然聊了三四个小时，而且聊得非常投入，这完全出乎他的意料。我看了托尼·罗宾斯书中精选出来的12位投资大师的采访记录，非常震惊，这些人，他能采访到就非常不容易了，能谈到这样深的程度，就更加不容易了。能把大师的知识融会贯通，总结提炼出简单7步投资法，让普通投资者即学即用，这是非常大的功德。只看大师的阵容名单，就会让你震惊：投资大师巴菲特、全球第一大对冲基金掌门人瑞·达利欧、全球最大指数基金公司先锋集团创始人约翰·博格、全球最大袭击并购投资人卡尔·伊坎、美国最大证券经纪公司嘉信理财创始人查尔斯·施瓦布，还有20世纪最伟大的投资人约翰·邓普顿爵士等。大师们的观点并不完全相同，但他们的基本原则是相同的。我们可以先从简单易行的投资方法学起，在有了一定的经验积累和财富积累后，再根据个人的特点，更多地模仿最适合自己的一两位大师。

好了，我们终于来到简单7步投资法的最后一步了，我们马上要揭开财富的最大秘密。

第七步，大财富。

前面六步，其实讲的是小财，不是大财。最后第七步告诉我们真正的大财富在哪里。

对于整个人类来说，最大的财富不在过去而在未来。巴菲特说现在和过去不同了，现在是科技的时代。全球股票市值最高的五大公司全部是科技公司，这五大公司的市值占美国5000家上市公司市值总和的1/10，而几十年前，市值最高的前五大公司都是银行和制造企业。对于我们每个人来说，最大的财富不是金钱，而是激情，感受到生命的激情，才是最大的财富。你活着不是为了钱，而是要用钱活得更好更幸福。这个主题是托尼·罗宾斯的特长，他研究了30多年，总结出做好三大决策关系到你人生的最大幸福。第一，关注什么。你关注的是你得到的还是你失去的，你能控制的还是你不能控制的。第二，有什么意义。意义就等于情绪，情绪就等于生命。有意义你才有激情，有激情才是真正的生命。第三，要做什么。关注带来意义，意义带来激情，激情激发状态，状态激发行动——持续的行动，行动最终带来改变，从此改变你的人生。

那么财富最终的秘密是什么？它是给予。我们讲的都是如何得到，但是能给你带来最大幸福的不是得到，不是得到更多，不是得到更好，而是给予。托尼·罗宾斯告诉我们，给不认识的陌生人餐食，我们的快乐程度还要再高出10倍。给予你不认识的陌生人，给你带来的幸福感最大。这被很多权威的心理学实验证实。托尼·罗宾斯说的最多的是他的个人经历。他儿时家徒四壁，一家人饿肚子，感恩节（相当于中国的春节）那一天，他推开门发现一篮子食物，是陌生人送来的，没有留下姓名。托尼·罗宾斯深受感动。后来他开始效仿这位好心的陌生人，给不认识的挨饿家庭提供餐食。从一年2家，到4家，再到8家，到现在他成立的基金会每年能够捐助5000万份餐食。我们都想得到最多的钱，最多的财富，可是我们没有想到，最大的财富不是得到而是给予，最大的财富不是过去而是未来，最大的财富不是物质而是激情。

无论是打球、跑步，还是写作、演讲，哪一项能力不是经历长时间练习才能培养出来的。这本书就是一本详细的理财训练手册，厚是必须的，因为行动需要详细具体的指导，需要重复练习来巩固提高。托尼·罗宾斯这个投资圈外人把投资理财讲得非常简单易行，让你愿意听、愿意学、愿意做。我真心推荐给你，我的朋友，这本书给你增加的财富不可估量。

活着只是为了钱，这没有意义，但是用钱来追求你的自由和梦想，这非常有意义。你做投资理财，不是为了钱，而是为了自己、为了家人、为了你要帮助的陌生人，你帮助他们生存、获得温饱、感到幸福、追求自由。

{\kaishu \hfill 刘建位

\hfill 2018年1月17日，于上海}
